\begin{problem}[第36届全国中学生物理竞赛决赛][光镊原理与实验分析]{115}
    2018年诺贝尔物理学奖的一部分授予美国科学家阿瑟·阿什金对光捕获（光镊）所做的开创性工作。光镊是用聚焦的激光束实现对微小颗粒的捕获，图a所示为一种光镊实验装置示意图。He-Ne激光器输出波长 $\lambda = 632.8\,\mathrm{nm}$ 的激光，光束直径 $D_{1} = 1.5\,\mathrm{mm}$，输出功率 $P = 30\,\mathrm{mW}$。为了使得捕获效果更好，可在输出的激光束后加一个扩束镜。激光束经反射镜反射后入射到一个焦距 $f = 4.0\,\mathrm{mm}$ 的聚焦物镜，聚焦物镜将光束汇聚于样品池内蒸馏水液体中，样品池内液体中有一定数量的聚苯乙烯小球，聚苯乙烯小球经过激光聚焦光斑时有可能被光捕获。已知小球的直径 $D_{\mathrm{b}} = 2.0\,\mu\mathrm{m}$，密度为 $\rho_{\mathrm{b}} = 1.0\times 10^{3}\,\mathrm{kg}\cdot \mathrm{m}^{-3}$，样品池光学玻璃材料折射率 $n = 1.46$，池内材料折射率 $n_{g}=1.46$，池内蒸馏水折射率 $n_{w}=1.33$，样品池底部玻璃厚度 $d=1.0\,\mathrm{mm}$，池底部外表面到聚焦物镜中心的距离 $a=2.0\,\mathrm{mm}$ 实验时温度 $T=300\,\mathrm{K}$，空气折射率 $n_{a}=1.0$，重力加速度 $g=9.80\,\mathrm{m\cdot s^{-2}}$，玻尔兹曼常量 $k_{B}=1.38\times10^{-23}\,\mathrm{J\cdot K^{-1}}$。
    \begin{subproblem}
        求不加扩束镜时激光束在空气中经聚焦物镜聚焦后聚焦光斑的直径；
    \end{subproblem}
    \begin{subproblem}
        为了提高捕获效果，对输出激光束进行扩束，采用的扩束镜由一个焦距 $f_{1} = -6.3\,\mathrm{mm}$ 的凹透镜和焦距 $f_{2} = 25.2\,\mathrm{mm}$ 的凸透镜组成，求两个透镜之间的距离和扩束后的平行激光束的直径 $D_{2}$；
    \end{subproblem}
    \begin{subproblem}
        将扩束后的激光束经图 a 所示聚焦物镜聚焦于样品池内蒸馏水中，求在傍轴近似下聚焦光斑中心到样品池底部内表面距离和聚焦光斑的直径；
    \end{subproblem}
    \begin{subproblem}
        图 b 为浸没在蒸馏水中的小球在梯度光强（简化成强度不同的两束光）中受力的几何光学模型，入射光线 $L_{1}$ 和 $L_{2}$ 沿 $x$ 轴方向，其光功率分别为 $p_{1}=1.0\,\mathrm{mW}$ 和 $p_{2}=2.0\,\mathrm{mW}$，光线经小球折射后其方向与 $x$ 轴的夹角如图 b 所示，不考虑光线在传播过程中的吸收和反射损耗，求小球在 $x$ 方向和 $y$ 方向所受到的力，并分别求其与小球重力之比；
    \end{subproblem}
    \begin{subproblem}
        求室温（$300\,\mathrm{K}$）下蒸馏水中小球的方均根速率；
    \end{subproblem}
    \begin{subproblem}
        光镊也可以通过光势阱模型理解，平行激光束经透镜聚焦后在径向 $\hat{r}$ 的光势阱分布为
        \begin{equation}
            U(r) = \begin{cases}
                -U_{0}\left(1 - \frac{r^{2}}{b^{2}}\right), & \text{当 } r \leq b \\
                0, & \text{当 } r > b
            \end{cases}
        \end{equation}
        上式中，$U_{0}=0.078\,\mathrm{eV}$，$b$ 是常量，$r$ 是到对称轴的距离。试问速率不超过多大的小球才能被光势阱捕获？
    \end{subproblem}
    \begin{subproblem}
        已知聚苯乙烯小球的速率分布函数可表示为：
        \begin{equation}
            f(v) = 4\pi \left(\frac{m}{2\pi k_{\mathrm{B}} T}\right)^{3/2} v^{2} e^{- \frac{m v^{2}}{2 k_{\mathrm{B}} T}}
        \end{equation}
        式中，$m$ 和 $v$ 分别是小球的质量和速率，$T$ 是绝对温度；试分析比较质量分别为 $m$ 和 $2m$ 的聚苯乙烯小球被光势阱捕获的概率大小。
    \end{subproblem}
\end{problem}